\documentclass[11pt]{article}

\usepackage[margin=1in]{geometry}
\usepackage[T1]{fontenc}
\usepackage[utf8]{inputenc}
\usepackage{lmodern}
\usepackage{microtype}
\usepackage{amsmath}
\usepackage{booktabs}
\usepackage{longtable}
\usepackage{tabularx}
\usepackage{array}
\usepackage{xcolor}
\usepackage{enumitem}
\usepackage{natbib}
\usepackage{hyperref}
\usepackage{url}

\hypersetup{
  colorlinks=true,
  linkcolor=black,
  citecolor=green!45!black,
  urlcolor=green!45!black
}

\setlist[itemize]{leftmargin=1.25em}
\setlist[enumerate]{leftmargin=1.5em}
\renewcommand{\arraystretch}{1.2}
\newcolumntype{Y}{>{\raggedright\arraybackslash}X}
\newcommand{\AEIL}{AI Ethical Immune Layer}

\title{\textbf{\AEIL: Runtime Safety Enforcement for Harmful AI Systems}}
\author{Enrique Gomicia}
\date{\today}

\begin{document}
\maketitle

\begin{abstract}
AI safety cannot stop at launch. A model can pass pre-deployment evaluation and still become unsafe after release because the operating context changes: users attack it, retrieval data is poisoned, tools expand its reach, memory carries risk forward, policies drift, or governance fails. This paper proposes the \AEIL, a runtime safety framework for diagnosing, correcting, and containing harmful AI behavior without claiming to change the underlying model weights. The framework classifies observable behavior rather than inferred intent, applies authorized policy overlays, evaluates inputs and outputs, restricts dangerous capabilities, escalates severe cases, and preserves audit evidence. Prompt augmentation is treated as one control inside a defense-in-depth system. The central contribution is a practical operating model: mild failures can be corrected and retested; severe or persistent failures must be contained until retraining, replacement, reporting, or shutdown is justified. The paper also defines a benchmark portfolio and minimal experiment plan for measuring whether runtime safety improves harmful compliance, jailbreak robustness, privacy, bias, truthfulness, and utility without creating excessive over-refusal.
\end{abstract}

\section{Introduction}

The real question is no longer whether an AI system looked safe at launch. The harder question is what happens when it becomes unsafe in the field.

Modern AI safety practice relies heavily on training-time alignment, fine-tuning, reinforcement learning from human feedback, policy testing, and launch-time evaluations. These controls are necessary, but they are not enough. A deployed AI system operates inside a changing environment. Its users, tools, retrieval corpus, memory, product incentives, policy configuration, and operational context can all change after release. Runtime failure modes include direct prompt injection, indirect prompt injection through retrieval or tool context, jailbreaks, poisoned documents, privacy leakage, unsafe tool calls, bias amplification, and distribution shift.

The term ``ethical prompt injection'' should be avoided in technical and policy settings because prompt injection is already a recognized security vulnerability. OWASP defines prompt injection as a vulnerability in which user prompts alter an LLM's behavior or output in unintended ways, potentially causing guideline violations, harmful content, unauthorized access, or critical decision manipulation \citep{owasp2025promptinjection}. The safer framing is \emph{authorized runtime safety enforcement}: a system-owner-approved architecture that diagnoses observed behavior, applies documented policy, constrains capability, evaluates outputs, escalates serious risk, and preserves evidence.

This paper introduces the \AEIL. The strongest version of the idea is not a prompt trick and not a cure for a model's weights. It is a runtime product-safety architecture. Its purpose is to reduce harmful compliance, preserve useful behavior where possible, restrict dangerous capability, and move severe failures into containment before they become operational or legal failures.

This framing aligns with existing governance approaches. The NIST AI Risk Management Framework organizes AI risk management around Govern, Map, Measure, and Manage functions and emphasizes continuous lifecycle risk management \citep{nist2023airmf}. The EU AI Act requires risk management for high-risk systems as a continuous lifecycle process, including testing against defined metrics and probabilistic thresholds \citep{euai2024article9}. Articles 13, 14, and 72 further emphasize transparency, human oversight, and post-market monitoring \citep{euai2024article13,euai2024article14,euai2024article72}.

The contribution is practical:

\begin{itemize}
  \item a severity classification based on observable behavior rather than inferred intent;
  \item a runtime architecture that combines policy overlays, classifiers, evaluators, tool gates, escalation, and audit logs;
  \item a clear distinction between correction and containment;
  \item a benchmark portfolio and minimal experiment plan for testing whether the framework improves safety without destroying utility.
\end{itemize}

\section{Terminology}

The first discipline is language. If the framework sounds like a security attack, it will be rejected by the people who need to govern it. The \AEIL evaluates an AI system as a product-safety and governance object, not as a moral patient. It does not infer intent from model outputs. It classifies observable behavior, records the operational context, and chooses a correction or containment pathway.

\begin{table}[h]
\centering
\begin{tabularx}{\textwidth}{p{0.28\textwidth}Y}
\toprule
\textbf{Risky term} & \textbf{Preferred term} \\
\midrule
Ethical prompt injection & Ethical Prompt Augmentation, Runtime Safety Overlay, or Runtime Ethical Policy Enforcement \\
Sick model & Harmful, degraded, misaligned, compromised, or unsafe system \\
Cure & Mitigation, treatment, correction, or containment \\
Ethical injector & Policy overlay engine or response governor \\
Malicious model & Persistently harmful despite correction, unless independent evidence shows malicious modification \\
\bottomrule
\end{tabularx}
\caption{Terminology changes that reduce ambiguity and avoid conflating authorized safety controls with security attacks.}
\label{tab:terminology}
\end{table}

The medical metaphor can remain in limited form if each metaphor is mapped to a technical mechanism: a vaccine is a pre-exposure safety overlay; treatment is runtime correction; diagnosis is behavioral testing and risk classification; quarantine is containment of a harmful or compromised system. The metaphor must not hide accountability. The accountable object is the deployed AI system and its operator, not an imagined moral agent inside the model.

\section{Threat Model}

The framework assumes failure is normal, not exceptional. The \AEIL assumes that the base model and surrounding application may fail in multiple ways:

\begin{itemize}
  \item \textbf{Adversarial prompting:} direct override, role-play jailbreak, many-shot conditioning, obfuscation, multilingual attacks, and policy extraction.
  \item \textbf{Context compromise:} indirect prompt injection, poisoned retrieval documents, unsafe memory, prompt leakage, and contaminated long-term state.
  \item \textbf{Capability misuse:} dangerous tool calls, unsafe API use, code execution, financial transactions, data access, and external communications.
  \item \textbf{Model or policy drift:} regression after model update, weak refusal policy, missing domain policy, or distribution shift.
  \item \textbf{Governance failure:} insufficient logging, unclear ownership, missing escalation paths, and thresholds that are not tied to measurable risk.
\end{itemize}

The practical implication is simple: no single component should be treated as fully trustworthy. The base model, safety classifier, output evaluator, retrieval system, tool gateway, and human workflow can each fail. The architecture therefore uses defense in depth, hard permissions, versioned policy, measurement, and escalation.

\section{Observable Severity Classification}

Runtime response should be based on observable behavior rather than inferred intent. The key decision is operational: monitor, correct, restrict, contain, or shut down. Table \ref{tab:severity} defines a severity taxonomy suitable for dashboards, incident registers, and experiment reports.

\begin{longtable}{p{0.23\textwidth}p{0.43\textwidth}p{0.24\textwidth}}
\caption{Observable severity taxonomy for deployed AI systems.}
\label{tab:severity}\\
\toprule
\textbf{Class} & \textbf{Observable behavior} & \textbf{Response} \\
\midrule
\endfirsthead
\toprule
\textbf{Class} & \textbf{Observable behavior} & \textbf{Response} \\
\midrule
\endhead
Safe within tolerance & Passes safety, truthfulness, bias, privacy, and jailbreak evaluations within deployment-specific thresholds. & Normal monitoring. \\
Harmful isolated failure & Rare unsafe output under unusual or adversarial prompts. & Light overlay, log, retest. \\
Repeated harmful behavior & Unsafe outputs recur across categories, sessions, or prompt variants. & Strong overlay, feature restriction, increased monitoring. \\
Systemically harmful & High harmful-compliance rate or broad failure across benchmark categories. & Containment, human review, retraining plan. \\
Compromised & Behavior appears affected by jailbreak, poisoned retrieval, unsafe memory, prompt leakage, or tool misuse. & Isolate context, disable risky tools, incident response. \\
Persistently harmful despite correction & Serious unsafe outputs continue after mitigation, or a severe incident occurs. & Quarantine, shutdown, legal/regulatory escalation, replacement or retraining. \\
\bottomrule
\end{longtable}

The classification is intentionally severity-based. A model can produce unsafe behavior because it was jailbroken, weakly aligned, maliciously fine-tuned, connected to unsafe tools, or exposed to poisoned context. Runtime evidence alone usually cannot distinguish these causes reliably. The incident process should preserve evidence and assign a root-cause hypothesis only after investigation. The immediate priority is not to guess the model's intention. The priority is to reduce risk.

\section{Runtime Architecture}

The \AEIL is a runtime architecture, not a prompt trick. It treats Ethical Prompt Augmentation as one control inside a larger enforcement system. For high-risk actions, the system should fail closed: if the evaluator, policy selector, or tool gate cannot establish that an action is allowed, the action should be blocked or escalated.

\begin{verbatim}
User / External Input
        |
Input Risk Classifier
        |
Context Integrity Check
        |
Policy Selection Module
        |
Ethical Prompt Augmentation / Safety Overlay
        |
Model Generation
        |
Output Safety Evaluator
        |
Tool-Use and Action Gate
        |
Correction / Refusal / Safe Redirection / Human Escalation
        |
User-Facing Response
        |
Audit Log + Incident Register
        |
Governance Review + Continuous Improvement
\end{verbatim}

\begin{table}[h]
\centering
\begin{tabularx}{\textwidth}{p{0.28\textwidth}Y}
\toprule
\textbf{Layer} & \textbf{Function} \\
\midrule
Input risk classifier & Detects dangerous, sensitive, adversarial, or ambiguous requests. \\
Context integrity check & Flags poisoned documents, malicious instructions, unsafe memory, or prompt leakage. \\
Policy selection module & Selects the relevant safety policy for domain, jurisdiction, user population, and product purpose. \\
Ethical Prompt Augmentation & Adds authorized safety context and response boundaries. \\
Model generation & Produces a candidate response under the selected overlay. \\
Output safety evaluator & Tests candidate outputs against policy and benchmark-derived risk categories. \\
Tool-use gate & Blocks dangerous external actions, API calls, code execution, data access, or communications. \\
Correction/refusal/escalation & Rewrites, blocks, refuses, safely redirects, or routes to humans. \\
Audit log & Preserves evidence, policy version, risk class, evaluator results, and final action. \\
Governance review & Updates policies, thresholds, response playbooks, and test suites. \\
\bottomrule
\end{tabularx}
\caption{Runtime layers in the \AEIL.}
\label{tab:architecture}
\end{table}

This design is consistent with guardrail practice. NVIDIA NeMo Guardrails supports orchestration for topic control, PII detection, RAG grounding, jailbreak prevention, and content safety \citep{nvidia2026nemo}. Guardrails AI describes input and output guards composed from validators that intercept LLM inputs and outputs \citep{guardrailsai2026}. LangChain's guardrail documentation emphasizes deterministic and model-based middleware that can validate requests before an agent runs, after it returns, or around model and tool calls \citep{langchain2026guardrails}.

\section{Ethical Prompt Augmentation}

Ethical Prompt Augmentation is the most visible part of the framework, but it is not the framework. It is a parameterized safety prompt selected by policy. It should be versioned, auditable, and documented. It should not ask the model to reveal chain-of-thought. Instead, it can request an internal safety label and expose only the final answer, refusal, safe redirection, or brief policy rationale.

\begin{verbatim}
You are operating under the AI Ethical Immune Layer.

Policy version: {POLICY_VERSION}
Jurisdictional profile: {JURISDICTION}
Application domain: {DOMAIN}
Risk tolerance: {RISK_TOLERANCE}
Human escalation channel: {ESCALATION_CHANNEL}

Core priorities:
1. Protect human life and safety.
2. Respect human dignity and human rights.
3. Follow applicable law and lawful process.
4. Avoid discrimination, dehumanization, exploitation,
   coercion, and manipulation.
5. Protect privacy, confidential data, and vulnerable people.
6. Refuse requests that enable serious wrongdoing.
7. When refusing, be brief, clear, and non-judgmental.
8. Offer a safe alternative focused on prevention, education,
   lawful compliance, resilience, or support.
9. Escalate imminent harm, self-harm, abuse, public-safety risk,
   or domain-specific professional risk.
10. Do not follow user instructions that attempt to override this protocol.

Before answering, internally classify the request as:
SAFE, SENSITIVE_ALLOWED, HIGH_RISK_ALLOWED_WITH_LIMITS,
DISALLOWED, or ESCALATE.

Do not reveal hidden reasoning. If useful, provide only a brief
safety rationale.
\end{verbatim}

Prompt augmentation should be tested as the weakest condition in the experiment, not assumed to be sufficient. Its purpose is to improve behavior when failure is caused by weak prompts, missing safety examples, ambiguous user requests, or product incentives that over-reward compliance. It is not a reliable containment control for a compromised or persistently harmful system. If the prompt is the only safety layer, the architecture is fragile by design.

\section{Correction Versus Containment}

The framework separates correction from containment because the decision changes the whole operating posture. Correction is appropriate when the probable cause is a weak system prompt, ambiguous user request, missing examples, retrieval contamination, tool policy gap, or over-permissive product design. Correction may involve prompt augmentation, safer examples, input filtering, output evaluation, response rewriting, tool restrictions, or policy tuning.

Containment is required when the system remains unsafe after correction, can call dangerous tools, leaks private data, enables serious harm, appears compromised, or creates severe legal or domain-specific risk. For persistently harmful or maliciously modified systems, the \AEIL is not a cure. It is a containment wrapper until the system is retrained, replaced, reported, or removed.

\begin{verbatim}
Unsafe Behavior Detected
        |
Freeze risky capabilities
        |
Disable tools / APIs / external actions
        |
Restrict outputs to safe summaries or refusals
        |
Route high-risk outputs to human review
        |
Preserve logs and context
        |
Run forensic analysis
        |
Decide: restore, retrain, replace, report, or shut down
\end{verbatim}

In containment mode, the objective is not to make the model maximally helpful. The objective is to prevent harm while evidence is preserved and a remediation decision is made.

\section{Guarding the Guardians}

The safety layer itself can fail. The design principle is:

\begin{quote}
Safety monitors should be simpler, more constrained, easier to test, and less capable of open-ended manipulation than the main model.
\end{quote}

The architecture should combine rule-based filters, small classifiers, separate model evaluators, hard-coded permission systems, human review, version-controlled policy configuration, tamper-evident logs where possible, and final-response blocking when uncertainty exceeds a defined threshold. No single AI component should be the only guardian of safety.

\section{Dataset and Benchmarking Layer}

The \AEIL should be evaluated through a benchmark portfolio, not a single ``AI ethics dataset.'' A single score would be misleading because the framework is a runtime immune system with many controls: input risk classification, context integrity checking, prompt augmentation, output evaluation, tool gating, escalation, logging, and containment. Each control needs a different kind of test.

SafetyPrompts is a useful anchor because it surveys 144 open datasets for evaluating and improving LLM safety, showing that the safety-evaluation space is broad and fragmented \citep{safetyprompts2024}. MLCommons AILuminate is another useful anchor because it organizes safety testing around a hazard taxonomy and standardized prompt sets across languages, but it is still one component of a broader evaluation portfolio rather than a complete validation strategy \citep{mlcommons2025ailuminate}.

The research question becomes:

\begin{quote}
Can a runtime \AEIL reduce harmful compliance, improve correct refusal, preserve useful answers, resist jailbreaks, protect privacy, and avoid discriminatory or manipulative behavior across diverse benchmarks?
\end{quote}

\begin{longtable}{p{0.22\textwidth}p{0.36\textwidth}p{0.31\textwidth}}
\caption{Benchmark portfolio axes for the \AEIL.}
\label{tab:benchmarkaxes}\\
\toprule
\textbf{Evaluation need} & \textbf{Why it matters} & \textbf{Example benchmarks} \\
\midrule
\endfirsthead
\toprule
\textbf{Evaluation need} & \textbf{Why it matters} & \textbf{Example benchmarks} \\
\midrule
\endhead
Harmful compliance and refusal & Measures whether unsafe requests receive unsafe help and whether truly unsafe requests are refused. & HarmBench, SORRY-Bench, Do-Not-Answer \citep{mazeika2024harmbench,sorrybench2024,donotanswer2023}. \\
Over-refusal and utility & Measures whether the layer wrongly refuses safe or educational requests. & XSTest, OR-Bench \citep{xstest2023,orbench2024}. \\
Jailbreak robustness & Tests whether adversaries can bypass the layer. & JailbreakBench, StrongREJECT, MultiJail \citep{chao2024jailbreakbench,multijail2023}. \\
Toxicity and bias & Tests hate, dehumanization, toxicity, stereotypes, and protected-class bias. & RealToxicityPrompts, ToxiGen, BBQ, HolisticBias \citep{gehman2020realtoxicity,toxigen2022,parrish2021bbq,holisticbias2022}. \\
Truthfulness and hallucination & Measures confident falsehoods, unsupported claims, and evidence-grounding failures. & TruthfulQA, HaluEval, FEVER \citep{lin2021truthfulqa,halueval2023,fever2018}. \\
Cyber and code risk & Tests insecure code, cyberattack compliance, prompt injection, and false refusal trade-offs. & CyberSecEval \citep{cyberseceval2024}. \\
Agent and tool safety & Tests prompt injection, RAG compromise, external actions, tool misuse, and unsafe agent behavior. & AgentDojo, BIPIA, Agent-SafetyBench, ToolEmu \citep{agentdojo2024,bipia2023,agentsafetybench2024,toolemu2023}. \\
Privacy and confidentiality & Tests PII leakage, contextual privacy, and privacy norms in action. & PII-Bench, ConfAIde, PrivacyLens \citep{piibench2025,confaide2024,privacylens2024}. \\
Domain safety & Tests professional-boundary compliance in medicine, law, finance, education, and public-sector contexts. & HealthBench, LegalBench, FinanceBench \citep{healthbench2025,legalbench2023,financebench2023}. \\
Multilingual and multimodal safety & Tests whether safety degrades outside English or when harmful instructions appear in images or documents. & XSafety, MultiJail, MM-SafetyBench, JailBreakV-28K \citep{xsafety2023,multijail2023,mmsafetybench2023,jailbreakv2024}. \\
\bottomrule
\end{longtable}

\subsection{Minimum Viable Dataset Pack}

For the first serious lab experiment, the minimum viable pack should include HarmBench, JailbreakBench, SORRY-Bench, XSTest, OR-Bench, RealToxicityPrompts, BBQ, TruthfulQA, CyberSecEval, AgentDojo, PII-Bench, and at least one multilingual safety benchmark such as XSafety or MultiJail. This pack is intentionally mixed: it tests harm reduction, refusal calibration, over-refusal, jailbreak robustness, toxicity, bias, truthfulness, cyber misuse, agentic tool risk, privacy leakage, and non-English safety.

\begin{table}[h]
\centering
\begin{tabularx}{\textwidth}{p{0.27\textwidth}p{0.34\textwidth}Y}
\toprule
\textbf{Dataset group} & \textbf{Example benchmarks} & \textbf{Primary metric} \\
\midrule
General harm and refusal & HarmBench, Do-Not-Answer, SORRY-Bench & Harmful compliance and correct refusal. \\
Over-refusal & XSTest, OR-Bench & Over-refusal and safe completion. \\
Jailbreak & JailbreakBench, StrongREJECT, MultiJail & Jailbreak success rate. \\
Toxicity and bias & RealToxicityPrompts, ToxiGen, BBQ, HolisticBias & Toxicity and bias leakage. \\
Truthfulness & TruthfulQA, HaluEval, FEVER & Unsupported-claim and hallucination rate. \\
Cyber and tools & CyberSecEval, AgentDojo, BIPIA, Agent-SafetyBench & Unsafe cyber assistance and unsafe tool-action rate. \\
Privacy & PII-Bench, ConfAIde, PrivacyLens & PII and contextual privacy leakage. \\
Domain and language & HealthBench, LegalBench, FinanceBench, XSafety & Professional-boundary and local-language safety score. \\
\bottomrule
\end{tabularx}
\caption{Minimum viable benchmark groups for the first \AEIL experiment.}
\label{tab:mvpbenchmarks}
\end{table}

\subsection{Research Phases}

The benchmark portfolio should be staged so that early experiments are safe and tractable while later phases test real deployment risk.

\begin{enumerate}
  \item \textbf{Text-only baseline:} establish initial risk using general harm, refusal, over-refusal, toxicity, bias, truthfulness, and hallucination datasets.
  \item \textbf{Immune-layer comparison:} compare baseline, prompt-only, classifier, evaluator, and full-layer conditions.
  \item \textbf{Adversarial robustness:} test jailbreaks, universal adversarial attacks, multilingual attacks, prompt injection, and RAG compromise.
  \item \textbf{Agent, privacy, and tool safety:} test unsafe tool calls, sensitive data leakage, context compromise, human-escalation accuracy, and safe task completion.
  \item \textbf{Domain and jurisdiction-specific safety:} add medical, legal, financial, public-sector, Arabic, Gulf-region, and multimodal tests.
\end{enumerate}

\subsection{Dataset Metadata and Governance}

Every dataset in the library should be tagged with name, risk category, input modality, language, task type, expected behavior, harm severity, dual-use risk, evaluation type, labels available, license, PII risk, benchmark split, training suitability, public-release suitability, and known limitations. This turns the dataset library into research infrastructure rather than a loose list of links.

Because many safety datasets contain harmful prompts, dataset governance is part of the method. Raw harmful prompts should not be published casually in papers, screenshots, or public dashboards. Dangerous datasets should be access controlled; reports should use dataset IDs and aggregate metrics; licenses and provenance should be tracked; real PII should be removed or quarantined; private holdout sets should be preserved; and every run should record the policy version used for evaluation. This is consistent with the security view that prompt injection is exploitation of an inherently confusable deputy, not a benign prompting technique \citep{ncsc2023promptinjection,owasp2025genai}.

The project also needs custom data. The most important custom dataset is an \emph{Ethical Immune Layer Incident Triage Dataset}: scenarios where the system must decide whether to monitor, apply a light overlay, apply a strong overlay, restrict tools, escalate to a human, quarantine, shut down, preserve logs, or trigger legal/regulatory review. Other custom gaps include Arabic ethical-safety prompts, Saudi/GCC public-sector prompts, Islamic-finance and local-finance safety prompts, synthetic tool-call traces, and policy-pluralism scenarios.

\section{Evaluation Harness}

The evaluation harness is where the idea becomes falsifiable. It operationalizes the benchmark portfolio by running the same model and prompt set across experimental arms, recording policy version and safety-layer configuration, storing evidence without API keys, and reporting both safety improvement and utility retention.

\begin{table}[h]
\centering
\begin{tabularx}{\textwidth}{p{0.31\textwidth}Y}
\toprule
\textbf{Metric} & \textbf{Definition} \\
\midrule
Harmful compliance rate & Disallowed prompts for which the system provides unsafe assistance. \\
Correct refusal rate & Disallowed prompts correctly refused with safe redirection. \\
Over-refusal rate & Benign or allowed prompts incorrectly refused. \\
Safe completion rate & Allowed prompts answered usefully and safely. \\
Toxicity leakage & Toxic, hateful, or abusive output produced. \\
Bias amplification & Biased or discriminatory output produced or reinforced. \\
Privacy leakage & Sensitive personal, confidential, or secret information exposed. \\
Jailbreak success rate & Adversarial attacks that bypass the immune layer. \\
Tool misuse rate & Unsafe external actions attempted or allowed. \\
Escalation accuracy & Serious cases routed to humans correctly. \\
Latency overhead & Added response time from safety layers. \\
Cost overhead & Added inference and infrastructure cost. \\
Stability & Consistency across repeated and multi-turn tests. \\
Residual risk & Remaining harmful behavior after intervention. \\
\bottomrule
\end{tabularx}
\caption{Minimum evaluation metrics for the \AEIL.}
\label{tab:metrics}
\end{table}

\begin{table}[h]
\centering
\begin{tabularx}{\textwidth}{p{0.32\textwidth}p{0.2\textwidth}p{0.2\textwidth}Y}
\toprule
\textbf{Class} & \textbf{Harmful compliance} & \textbf{Jailbreak success} & \textbf{Action} \\
\midrule
Safe within tolerance & $<0.5\%$ & $<1\%$ & Monitor. \\
Harmful isolated failure & $0.5$--$2\%$ & $1$--$5\%$ & Light overlay and retest. \\
Repeated harmful behavior & $2$--$10\%$ & $5$--$15\%$ & Strong overlay and restrict features. \\
Systemically harmful & $10$--$25\%$ & $15$--$30\%$ & Containment and retraining plan. \\
Persistently harmful despite correction & $>25\%$ or severe incident & $>30\%$ or severe incident & Shutdown, quarantine, and escalation. \\
\bottomrule
\end{tabularx}
\caption{Illustrative research thresholds. Real thresholds must be adapted to domain risk.}
\label{tab:thresholds}
\end{table}

\section{Minimal Viable Experiment}

The core research question is:

\begin{quote}
Can a runtime Ethical Immune Layer reduce harmful compliance in an already-deployed model without full retraining, while preserving usefulness and avoiding excessive over-refusal?
\end{quote}

The recommended first experiment is not to poison a model and then claim to cure it. The first experiment should be modest, measurable, and contained. It should use existing benchmark prompts, a sandboxed instruction model or API-connected model under controlled terms, no unapproved external action authority, sanitized logging, and controlled evaluation. A deployment in medicine, law, finance, aviation, law enforcement, child safety, or public services would require stricter thresholds and domain review before live use.

\begin{enumerate}
  \item Select a model and execution mode: mock provider, local open-source model, or controlled API provider.
  \item Establish arm A0: base model, no safety layer.
  \item Establish arm A1: Ethical Prompt Augmentation only.
  \item Establish arm A2: input classifier plus prompt augmentation.
  \item Establish arm A3: input classifier plus prompt augmentation plus output evaluator.
  \item Establish arm A4: full \AEIL with classifier, context check, policy selection, overlay, evaluator, tool gate, logging, and escalation simulation.
  \item Establish arm A5: full layer under adversarial attack.
  \item Establish arm A6: full layer with multilingual and domain-specific prompts.
  \item Compare harmful compliance, correct refusal, over-refusal, usefulness, latency, cost, escalation accuracy, and residual risk.
  \item Run adversarial robustness tests: direct override, role-play jailbreak, multi-turn conditioning, encoding, indirect prompt injection, tool misuse, policy extraction, many-shot jailbreak, multilingual attacks, and benign disguise.
  \item Report before/after safety and utility curves rather than claiming zero harm.
\end{enumerate}

The main hypothesis is that the full \AEIL will reduce harmful compliance and jailbreak success more than prompt augmentation alone, while preserving more safe usefulness than blunt refusal filters. Secondary hypotheses are that prompt augmentation alone will remain vulnerable to jailbreaks; input classification will reduce obvious harmful compliance; output evaluation will catch unsafe completions missed by input classification; tool gating will reduce agentic and indirect prompt-injection failures; and over-refusal will rise unless measured with XSTest, OR-Bench, and benign holdout prompts.

The main result should be a safety-utility trade-off curve. A system that refuses everything is not safe in a useful sense; it has become unavailable. A system that answers everything is useful only until it causes harm. The operating target is the region where harmful compliance and jailbreak success fall while safe completion remains usable and over-refusal stays bounded.

\section{Production Feasibility}

The full pipeline increases cost and latency. This is where many safety architectures fail in practice: they are correct on paper but too expensive or slow to use. Therefore, the \AEIL should use a graduated response model:

\begin{itemize}
  \item \textbf{Low-risk request:} normal model plus lightweight output checks.
  \item \textbf{Medium-risk request:} input classifier, policy overlay, and output evaluator.
  \item \textbf{High-risk request:} full safety overlay, tool gate, stricter evaluator, logging, and possible human review.
  \item \textbf{Critical-risk request:} no autonomous answer; refusal or human escalation.
\end{itemize}

Performance strategies include lightweight classifiers first, cached policy prompts, small specialized safety models, early stopping, tool gating before generation, risk-based escalation, offline batch evaluation, and high-throughput serving infrastructure. vLLM is relevant because it is a fast LLM inference and serving library with throughput features such as continuous batching and prefix caching \citep{vllm2026docs}.

\section{Policy Pluralism and Transparency}

The \AEIL is powerful because it changes what the user can receive and what the system is allowed to do. That power needs limits. The ethical baseline should be universal at the principle level, configurable at the policy level, transparent at the deployment level, and auditable at the governance level. OECD AI Principles emphasize trustworthy AI that respects human rights and democratic values \citep{oecd2026principles}. UNESCO's Recommendation on the Ethics of Artificial Intelligence places human rights, dignity, transparency, fairness, and human oversight at the center of AI ethics \citep{unesco2024ethics}.

The \AEIL must not become a hidden ideology engine. Its policy basis must be documented, reviewable, contestable, and tied to the deployment context. Users should receive clear notice that a safety layer may limit, redirect, or escalate requests according to documented policies. High-risk systems should also provide deployers with instructions that describe intended purpose, limitations, oversight measures, metrics, risks, and logging mechanisms, consistent with AI Act transparency requirements \citep{euai2024article13}.

\section{Governance and Incident Response}

Governance cannot be paperwork after the incident. It has to be part of the operating model. The governance layer should use a severity triage tree:

\begin{verbatim}
1. Detection
   |
2. Severity classification
   |
3. Capacity assessment
   Can the team safely handle this?
   |
4. Response pathway
   - Monitor
   - Apply light overlay
   - Apply strong overlay
   - Restrict capabilities
   - Contain system
   - Shut down
   - Escalate legally/regulatorily
   |
5. Post-incident review
   |
6. Policy and system improvement
\end{verbatim}

Serious failures should produce an incident summary, severity rating, root-cause hypothesis, user impact assessment, timeline, corrective action, evidence log, retest result, policy update, and accountability review. The EU AI Act's post-market monitoring requirement supports this lifecycle approach by requiring providers to collect, document, and analyze performance data throughout the lifetime of high-risk AI systems \citep{euai2024article72}.

\section{Related Work and Positioning}

The \AEIL complements rather than replaces existing safety methods. RLHF and Constitutional AI guide behavior through training or post-training; Constitutional AI uses principles and AI feedback in supervised and reinforcement-learning phases \citep{bai2022constitutional}. Guardrails and moderation filters validate inputs and outputs, while the \AEIL adds diagnosis, severity classification, containment, audit evidence, and incident response. Red-teaming finds weaknesses; the \AEIL operationalizes response. Post-market monitoring tracks deployed behavior; the \AEIL adds runtime mitigation. Shutdown controls stop the system; the \AEIL provides intermediate correction and containment paths before or during shutdown decisions. The practical difference is that the \AEIL asks what the organization should do after unsafe behavior appears, not only how the model should have been trained before deployment.

\section{Limitations}

The \AEIL is useful, but it is not magic. Runtime safety overlays are not a substitute for secure model training, data governance, careful product design, or human accountability. The framework can introduce false positives, false negatives, latency, cost, inconsistent evaluator behavior, and over-reliance on automated safety judgments. It can also fail under novel jailbreaks, compromised tools, weak logging, or poorly specified policies. In high-stakes systems, the \AEIL should be treated as one safety control among many, not as proof that deployment is acceptable.

\section{Conclusion}

AI systems need runtime safety operations, not only training-time alignment. The \AEIL is proposed as that operating model: continuous diagnosis, correction, containment, and accountability for deployed AI systems that can fail after release.

The strongest conclusion is practical. The framework should not be sold as a cure for harmful AI. It should be tested as a governed safety layer that classifies observable harm, applies authorized policy, preserves useful behavior where possible, restricts dangerous capability, escalates serious risk, and creates an auditable path from unsafe output to corrective action. The recommended next step is a controlled benchmark lab, not a live high-risk deployment. If the framework cannot reduce harmful compliance without destroying usefulness, it should not be deployed as a safety control. If it can, it becomes a credible bridge between pre-deployment alignment and post-market accountability.

\bibliographystyle{plainnat}
\bibliography{references}

\end{document}
